\documentclass[11pt,a4paper]{article}

% ============================================================
% PAQUETES
% ============================================================

\usepackage[utf8]{inputenc}
\usepackage[T1]{fontenc}
\usepackage[spanish,es-nodecimaldot]{babel}
\usepackage{lmodern}
\usepackage{textcomp}
\usepackage[margin=2.5cm]{geometry}
\usepackage{graphicx}
\usepackage{float}
\usepackage{xcolor}
\usepackage{tcolorbox}
\tcbuselibrary{skins,breakable}
\usepackage{booktabs}
\usepackage{longtable}
\usepackage{array}
\usepackage{tabularx}
\usepackage{colortbl}
\usepackage{amssymb}
\usepackage{pifont}
\usepackage{enumitem}
\usepackage{fancyhdr}
\usepackage{titlesec}
\usepackage{hyperref}

% ============================================================
% COLORES
% ============================================================

\definecolor{bancoblue}{HTML}{1A237E}
\definecolor{bancolight}{HTML}{E8EAF6}
\definecolor{bancored}{HTML}{C62828}
\definecolor{bancoorange}{HTML}{E65100}
\definecolor{bancogreen}{HTML}{2E7D32}
\definecolor{bancogray}{HTML}{4C4C4C}

% ============================================================
% HYPERREF
% ============================================================

\hypersetup{
	colorlinks=true,
	linkcolor=bancoblue,
	urlcolor=bancoblue,
	citecolor=bancoblue,
	pdftitle={Documento de Alcance del Proyecto - Banco del Lago, S.A.},
	pdfauthor={Walter Gómez}
}

% ============================================================
% ENCABEZADOS
% ============================================================

\pagestyle{fancy}
\fancyhf{}
\fancyhead[L]{\textcolor{bancogray}{\small Documento de Alcance --- Banco del Lago, S.A.}}
\fancyhead[R]{\textcolor{bancogray}{\small \today}}
\fancyfoot[C]{\thepage}
\fancyfoot[L]{\textcolor{bancogray}{\small Confidencial --- Uso interno}}
\renewcommand{\headrulewidth}{0.4pt}
\renewcommand{\footrulewidth}{0.4pt}

% ============================================================
% TÍTULOS
% ============================================================

\titleformat{\section}{\Large\bfseries\color{bancoblue}}{\thesection}{1em}{}
\titleformat{\subsection}{\large\bfseries\color{bancoblue}}{\thesubsection}{1em}{}
\titleformat{\subsubsection}{\normalsize\bfseries\color{bancogray}}{\thesubsubsection}{1em}{}

% ============================================================
% CAJAS
% ============================================================

\newtcolorbox{infobox}[1][]{
	colback=bancolight,
	colframe=bancoblue,
	fonttitle=\bfseries,
	title=#1,
	boxrule=0.5pt,
	arc=2pt,
	breakable
}

\newtcolorbox{warnbox}[1][]{
	colback=red!5,
	colframe=bancored,
	fonttitle=\bfseries,
	title=#1,
	boxrule=0.5pt,
	arc=2pt,
	breakable
}

% ============================================================
% SÍMBOLOS
% ============================================================

\newcommand{\checkmarkverde}{\textcolor{bancogreen}{\ding{51}}}
\newcommand{\cruzroja}{\textcolor{bancored}{\ding{55}}}

% ============================================================
% TÍTULO
% ============================================================

\title{
	\vspace{-1.5cm}
	\textbf{\Huge Documento de Alcance del Proyecto}\\[0.4cm]
	\Large Gap Analysis de Ciberseguridad y Riesgo Tecnológico\\[0.3cm]
	\large Banco del Lago, S.A.\\[0.5cm]
	\normalsize Project Charter / Acta de Constitución\\[0.3cm]
	\normalsize Versión 1.1 --- Septiembre 2026\\[0.5cm]
	\normalsize Preparado por: Walter Gómez\\
	\normalsize ISO 27001 Lead Implementer | DevSecOps Engineer
}
\date{}

% ============================================================
% DOCUMENTO
% ============================================================

\begin{document}
	
	\maketitle
	\thispagestyle{empty}
	
	\begin{infobox}[Control del documento]
		\textbf{Tipo de documento:} Project Charter / Acta de Constitución \\
		\textbf{Estado:} Aprobado para inicio de fase de diagnóstico \\
		\textbf{Versión:} 1.1 \\
		\textbf{Fecha:} Septiembre 2026 \\
		\textbf{Confidencialidad:} Uso interno --- caso de estudio con fines de portafolio profesional
	\end{infobox}
	
	\vspace{0.5cm}
	
	\tableofcontents
	
	\newpage
	
	% ============================================================
	% SECCIÓN 1
	% ============================================================
	
	\section{Propósito del documento}
	
	Este documento formaliza el alcance, los supuestos y las bases regulatorias del proyecto de \textbf{Gap Analysis de Ciberseguridad y Riesgo Tecnológico} para Banco del Lago, S.A., banco ficticio construido con datos reales del sistema financiero guatemalteco para fines de caso de estudio. Sirve como punto de partida único para todas las fases posteriores del proyecto (Gap Analysis, plan de remediación, roadmap de implementación).
	
	Todo dato en este documento está etiquetado según su origen:
	
	\begin{itemize}[leftmargin=1.5em]
		\item \textbf{[VERIFICADO]} --- dato real, con fuente y cálculo trazable.
		\item \textbf{[ESTIMADO]} --- derivado por proporción o razón a partir de un dato real.
		\item \textbf{[SUPUESTO]} --- decisión de diseño del caso de estudio, sin dato real detrás, pero técnicamente razonable.
	\end{itemize}
	
	% ============================================================
	% SECCIÓN 2
	% ============================================================
	
	\section{Perfil institucional}
	
	\subsection{Identidad}
	
	\begin{table}[H]
		\centering
		\begin{tabular}{@{}p{4cm}p{6cm}p{3cm}@{}}
			\toprule
			\textbf{Campo} & \textbf{Valor} & \textbf{Estatus} \\ \midrule
			Razón social & Banco del Lago, S.A. & Ficticio \\
			Grupo financiero & Grupo Financiero del Lago & Supuesto \\
			Sede & Ciudad de Guatemala & Supuesto \\
			Ente supervisor & Superintendencia de Bancos (SIB) & Verificado \\ \bottomrule
		\end{tabular}
	\end{table}
	
	\subsection{Tamaño y posición de mercado}
	
	Base de cálculo: Banco Promerica de Guatemala +15\%.
	
	\begin{longtable}{@{}p{6.5cm}p{5cm}p{3.5cm}@{}}
		\toprule
		\textbf{Campo} & \textbf{Valor} & \textbf{Estatus} \\ \midrule
		\endfirsthead
		\toprule
		\textbf{Campo} & \textbf{Valor} & \textbf{Estatus} \\ \midrule
		\endhead
		Empleados & $\sim$2,800 & Estimado \\
		Activos totales & $\sim$Q39,290 millones & Verificado \\
		Cartera de préstamos & $\sim$Q29,190 millones & Verificado \\
		Patrimonio & $\sim$Q4,223 millones & Verificado \\
		Índice de adecuación de capital & $\sim$13.5--14\% & Estimado \\
		Participación en activos del sistema & $\sim$6\% & Estimado \\
		Posición en el ranking bancario & 6\textdegree \, a 10\textdegree \, lugar & Estimado \\ \bottomrule
	\end{longtable}
	
	\subsection{Red de distribución y clientes}
	
	\begin{longtable}{@{}p{6cm}p{5cm}p{3cm}@{}}
		\toprule
		\textbf{Campo} & \textbf{Valor} & \textbf{Estatus} \\ \midrule
		\endfirsthead
		\toprule
		\textbf{Campo} & \textbf{Valor} & \textbf{Estatus} \\ \midrule
		\endhead
		Agencias & 85 & Supuesto \\
		Cajeros automáticos propios & 180 & Supuesto \\
		Clientes activos & $\sim$750,000 & Supuesto \\
		Canales digitales & Banca en línea y móvil & Supuesto \\
		Volumen mensual de transacciones & $\sim$2.8 millones & Supuesto \\
		Distribución por canal & 70\% digital / 30\% agencia & Supuesto \\ \bottomrule
	\end{longtable}
	
	\subsection{Infraestructura tecnológica}
	
	\begin{longtable}{@{}p{5cm}p{8.5cm}@{}}
		\toprule
		\textbf{Campo} & \textbf{Valor} \\ \midrule
		\endfirsthead
		\toprule
		\textbf{Campo} & \textbf{Valor} \\ \midrule
		\endhead
		Centro de cómputo principal & Ciudad de Guatemala \\
		Centro de cómputo alterno / DRP & Fuera del área metropolitana \\
		Core bancario & Producto comercial de terceros \\
		Modelo de nube & Híbrido (cargas no críticas en nube; core y datos sensibles on-premise) \\
		SOC / CSIRT & Tercerizado (MDR), sin capacidad 24/7 interna \\ \bottomrule
	\end{longtable}
	
	\subsubsection*{Proveedores tecnológicos críticos}
	
	\begin{longtable}{@{}cp{10cm}@{}}
		\toprule
		\textbf{\#} & \textbf{Proveedor} \\ \midrule
		\endfirsthead
		\toprule
		\textbf{\#} & \textbf{Proveedor} \\ \midrule
		\endhead
		1 & Proveedor de core bancario \\
		2 & Procesador/switch de tarjetas \\
		3 & Proveedor de nube \\
		4 & Proveedor de canales electrónicos / banca móvil \\
		5 & Proveedor de mensajería SWIFT \\
		6 & Proveedor de centro de cómputo alterno \\ \bottomrule
	\end{longtable}
	
	\subsection{Gobernanza}
	
	\begin{longtable}{@{}p{6.5cm}p{6cm}p{2cm}@{}}
		\toprule
		\textbf{Órgano / Rol} & \textbf{Situación} & \textbf{Estatus} \\ \midrule
		\endfirsthead
		\toprule
		\textbf{Órgano / Rol} & \textbf{Situación} & \textbf{Estatus} \\ \midrule
		\endhead
		Consejo de Administración & Responsable del riesgo tecnológico & Verificado \\
		Comité de Gestión de Riesgos & Existe y opera & Verificado \\
		Comité de Auditoría & Existe y opera & Verificado \\
		Unidad de Administración de Riesgos & Existe y opera & Verificado \\
		Auditoría Interna & Reporta al Comité de Auditoría & Verificado \\
		Oficial de Seguridad (CISO) & No existe formalmente & Supuesto \\ \bottomrule
	\end{longtable}
	
	% ============================================================
	% SECCIÓN 3
	% ============================================================
	
	\section{Marco regulatorio aplicable}
	
	\begin{longtable}{@{}p{2.5cm}p{8cm}p{4.5cm}@{}}
		\toprule
		\textbf{Norma} & \textbf{Descripción} & \textbf{Relevancia} \\ \midrule
		\endfirsthead
		\toprule
		\textbf{Norma} & \textbf{Descripción} & \textbf{Relevancia} \\ \midrule
		\endhead
		JM-98-2025 & Reglamento para la Administración del Riesgo Tecnológico. Vigente desde el 30-oct-2025. Deroga la JM-104-2021. & Marco normativo principal \\ \midrule
		JM-91-2024 / JM-99-2025 & Reglamento de Medidas de Seguridad en Canales Electrónicos y su modificación (autenticación reforzada). & Aplica a banca en línea y móvil \\ \midrule
		JM-62-2016 & Reglamento de Gobierno Corporativo. & Marco de gobernanza \\ \midrule
		Decreto 15-2026 & Ley Integral para la Prevención y Represión del Lavado de Dinero. Vigente desde el 17-sep-2026. & Contexto regulatorio ampliado \\ \midrule
		ISO/IEC 27001:2022 & Estándar internacional de gestión de seguridad de la información. No obligatorio en Guatemala. & Marco de referencia voluntario \\ \midrule
		PCI DSS & Estándar de la industria de tarjetas de pago. & Aplica por operación de tarjetas \\ \bottomrule
	\end{longtable}
	
	\subsection{Calendario de cumplimiento regulatorio}
	
	\begin{longtable}{@{}p{7.5cm}p{3.5cm}p{4cm}@{}}
		\toprule
		\textbf{Hito regulatorio} & \textbf{Fecha límite} & \textbf{Estatus} \\ \midrule
		\endfirsthead
		\toprule
		\textbf{Hito regulatorio} & \textbf{Fecha límite} & \textbf{Estatus} \\ \midrule
		\endhead
		Actualización Manual de Riesgo Tecnológico & Octubre 2026 & Pendiente verificación \\
		Actualización del DRP & Octubre 2026 & Pendiente verificación \\
		Clasificación de proveedores (Art. 52) & Dic 2026 -- Mar 2027 & Pendiente confirmación \\
		Revisión de canales electrónicos & Dic 2026 -- Mar 2027 & Estimado \\ \bottomrule
	\end{longtable}
	
	\begin{warnbox}[Implicación para el proyecto]
		El Gap Analysis debe completarse antes de octubre de 2026 para dar tiempo a la remediación de los plazos de octubre. La clasificación de proveedores debe estar lista antes de diciembre de 2026. La revisión de canales electrónicos debe coincidir con la próxima actualización regulatoria.
	\end{warnbox}
	
	% ============================================================
	% SECCIÓN 4
	% ============================================================
	
	\section{Incidente disparador del proyecto}
	
	Escenario compuesto por dos elementos documentados a nivel sectorial:
	
	\begin{enumerate}[leftmargin=1.5em]
		\item \textbf{Presión regulatoria con fecha límite.} El Manual de Administración del Riesgo Tecnológico y el Plan de Recuperación ante Desastres deben actualizarse dentro de los 12 meses posteriores a la entrada en vigor de la JM-98-2025 (octubre 2026). Adicionalmente, la clasificación de criticidad de proveedores tecnológicos debe completarse entre diciembre 2026 y marzo 2027, y se anticipa una nueva revisión de la regulación de canales electrónicos en el mismo período.
		\item \textbf{Patrón de fraude sectorial.} El Ministerio Público acumuló más de 56,000 denuncias por fraude electrónico en tres años a nivel nacional, con modalidades de suplantación de identidad y enlaces fraudulentos distribuidos por WhatsApp como las más frecuentes. Banco del Lago detecta un incidente de este tipo que expone una debilidad de autenticación en uno de sus canales electrónicos.
	\end{enumerate}
	
	La combinación de ambos elementos da al Comité de Gestión de Riesgos motivación suficiente para solicitar al Consejo de Administración la autorización del proyecto.
	
	% ============================================================
	% SECCIÓN 5
	% ============================================================
	
	\section{Objetivo del proyecto}
	
	Identificar las brechas de cumplimiento de Banco del Lago frente al Reglamento para la Administración del Riesgo Tecnológico (JM-98-2025) y frente a las buenas prácticas de ISO/IEC 27001:2022, y establecer una hoja de ruta priorizada de remediación que permita:
	
	\begin{itemize}[leftmargin=1.5em]
		\item Cumplir los plazos regulatorios de la JM-98-2025 (octubre 2026, diciembre 2026--marzo 2027).
		\item Formalizar la función de seguridad de la información (CISO / Oficial de Seguridad de la Información).
		\item Reducir la exposición al fraude en canales electrónicos.
		\item Preparar al banco para la próxima revisión regulatoria de canales electrónicos (JM-91-2024 / JM-99-2025).
	\end{itemize}
	
	% ============================================================
	% SECCIÓN 6
	% ============================================================
	
	\section{Alcance}
	
	\subsection{Incluido}
	
	\begin{itemize}[leftmargin=1.5em]
		\item Gobernanza de riesgo tecnológico (Consejo, Comité de Gestión de Riesgos, Unidad de Riesgos).
		\item Gestión de infraestructura de TI, sistemas de información y bases de datos.
		\item Seguridad de la información (control de accesos, cifrado, gestión de vulnerabilidades, pruebas de penetración).
		\item Capacidad de detección y respuesta (monitoreo continuo, SIEM, SOC tercerizado, gestión de incidentes).
		\item Continuidad de operaciones de TI (DRP, centro de cómputo alterno).
		\item Gestión y clasificación de criticidad de proveedores tecnológicos.
		\item Seguridad en canales electrónicos (banca en línea, banca móvil).
		\item Mapeo de controles contra ISO/IEC 27001:2022 como marco de referencia complementario.
	\end{itemize}
	
	\subsection{Excluido}
	
	\begin{itemize}[leftmargin=1.5em]
		\item Cumplimiento antilavado bajo el Decreto 15-2026 (se trata como contexto, no como objeto de auditoría).
		\item Certificación formal ISO/IEC 27001 (el proyecto usa la norma como marco de referencia, no busca la certificación).
		\item Auditoría de riesgo de crédito, mercado o liquidez.
		\item Evaluación de las demás empresas del Grupo Financiero del Lago fuera del banco.
	\end{itemize}
	
	% ============================================================
	% SECCIÓN 7
	% ============================================================
	
	\section{Gobernanza del proyecto}
	
	\begin{longtable}{@{}p{3.5cm}p{4.25cm}p{6cm}@{}}
		\toprule
		\textbf{Rol} & \textbf{Responsable} & \textbf{Función} \\ \midrule
		\endfirsthead
		\toprule
		\textbf{Rol} & \textbf{Responsable} & \textbf{Función} \\ \midrule
		\endhead
		Sponsor & Gerente General & Autoriza el proyecto y presenta resultados al Consejo \\
		Instancia de respaldo & Comité de Gestión de Riesgos & Da seguimiento y aprueba el plan de remediación \\
		Aprobación final & Consejo de Administración & Aprueba el resultado y el presupuesto \\
		Dueño operativo & Oficial de Seguridad (a formalizar) & Ejecuta el diagnóstico y coordina la remediación \\
		Revisión independiente & Auditoría Interna & Revisa el proceso y los resultados \\ \bottomrule
	\end{longtable}
	
	% ============================================================
	% SECCIÓN 8
	% ============================================================
	
	\section{Entregables}
	
	\begin{enumerate}[leftmargin=1.5em]
		\item Gap Analysis documentado contra JM-98-2025 e ISO/IEC 27001:2022 (matriz de cumplimiento por artículo/control).
		\item Registro de hallazgos priorizados por nivel de riesgo.
		\item Plan de remediación con responsables, plazos y dependencias regulatorias.
		\item Propuesta de estructura para la función de seguridad de la información.
		\item Dashboard de madurez con scoring por dominio del NIST CSF 2.0.
	\end{enumerate}
	
	% ============================================================
	% SECCIÓN 9
	% ============================================================
	
	\section{Criterios de éxito del proyecto}
	
	\begin{longtable}{@{}p{6cm}p{6cm}p{3cm}@{}}
		\toprule
		\textbf{Criterio} & \textbf{Métrica} & \textbf{Plazo} \\ \midrule
		\endfirsthead
		\toprule
		\textbf{Criterio} & \textbf{Métrica} & \textbf{Plazo} \\ \midrule
		\endhead
		Gap Analysis completado & Aprobado por el Comité de Riesgos & 4 semanas \\
		Plan de remediación aprobado & Con responsables y plazos & 6 semanas \\
		Propuesta de estructura CISO & Aprobada por el Consejo & 8 semanas \\
		Quick wins identificados & Al menos 3 en 90 días & 4 semanas \\
		Cumplimiento plazos regulatorios & Manual y DRP actualizados & Octubre 2026 \\
		Clasificación de proveedores & 6 proveedores críticos clasificados & Diciembre 2026 \\ \bottomrule
	\end{longtable}
	
	% ============================================================
	% SECCIÓN 10
	% ============================================================
	
	\section{Riesgos del proyecto}
	
	\begin{longtable}{@{}p{6.2cm}p{1cm}p{1.25cm}p{6cm}@{}}
		\toprule
		\textbf{Riesgo} & \textbf{Prob.} & \textbf{Impacto} & \textbf{Mitigación} \\ \midrule
		\endfirsthead
		\toprule
		\textbf{Riesgo} & \textbf{Prob.} & \textbf{Impacto} & \textbf{Mitigación} \\ \midrule
		\endhead
		Resistencia al cambio en TI & Alta & Medio & Involucrar a TI desde el inicio \\
		Falta de documentación existente & Alta & Alto & Entrevistas estructuradas; reconstruir inventario \\
		Presupuesto limitado & Media & Alto & Priorizar quick wins de bajo costo \\
		Dependencia de proveedores externos & Media & Medio & Cláusulas de seguridad; evaluar alternativas \\
		Plazos regulatorios ajustados & Alta & Alto & Completar Gap Analysis antes de octubre 2026 \\
		Falta de personal especializado & Media & Medio & Capacitar al equipo; considerar MDR \\ \bottomrule
	\end{longtable}
	
	% ============================================================
	% SECCIÓN 11
	% ============================================================
	
	\section{Próxima fase}
	
	Inicio del \textbf{Gap Analysis} contra JM-98-2025 e ISO/IEC 27001:2022, usando este documento como línea base del perfil institucional. La evaluación se estructurará contra los 6 dominios del NIST CSF 2.0 (Gobernar, Identificar, Proteger, Detectar, Responder, Recuperar), con scoring de madurez de 5 niveles.
	
	% ============================================================
	% SECCIÓN 12
	% ============================================================
	
	\section{Control de supuestos}
	
	Toda cifra marcada como \textbf{[SUPUESTO]} en este documento es una decisión de diseño del caso de estudio, no un dato de Banco Promerica de Guatemala ni de ningún otro banco real. Las cifras marcadas \textbf{[VERIFICADO]} provienen de fuentes públicas sobre Banco Promerica de Guatemala (informes financieros y de mercado a 2025--2026) con un ajuste del +15\% aplicado de forma explícita y trazable.
	
	\vspace{1cm}
	
	\begin{center}
		\textcolor{bancogray}{\rule{0.5\textwidth}{0.4pt}}\\[0.3cm]
		\textit{Documento preparado como parte del portafolio técnico de Walter Gómez.\\
			Proyecto 4: Modelo de Madurez.}\\[0.3cm]
		\textcolor{bancogray}{\rule{0.5\textwidth}{0.4pt}}
	\end{center}
	
\end{document}